\documentclass[12pt]{article}

\usepackage[a4paper,margin=1in]{geometry}
\usepackage{amsmath,amssymb}
\usepackage{tikz}
\usepackage{enumitem}

\setlength{\parindent}{0pt}
\setlength{\parskip}{0.75em}

\title{Lecture 11 Supplement: Conditional Poisson, Binomial Splitting, and Event Gaps}
\author{MA4034 -- Time Series and Stochastic Processes}
\date{}

\begin{document}

\maketitle

\section{Main Idea}

There are two common types of Poisson process questions.

\begin{enumerate}
    \item Questions about the \textbf{number of events in intervals}.
    \item Questions about the \textbf{time gap between events}.
\end{enumerate}

For counting events, we often use the Poisson distribution.

For waiting times between events, we often use the exponential distribution.

Sometimes, after we are told the total number of events, the problem becomes a binomial problem.

\section{Why Does the Binomial Distribution Appear?}

Suppose cars pass a point according to two independent Poisson processes:

\[
\text{Northbound cars: rate } 3 \text{ per minute},
\]

\[
\text{Southbound cars: rate } 2 \text{ per minute}.
\]

The total rate is
\[
3+2=5.
\]

So if we ignore direction, total cars pass according to a Poisson process with rate \(5\) per minute.

\section{Given That a Car Passed, What Is the Probability It Is Northbound?}

Since northbound cars have rate \(3\), and total cars have rate \(5\),

\[
P(\text{northbound})=\frac{3}{5}.
\]

Similarly,

\[
P(\text{southbound})=\frac{2}{5}.
\]

This is the key step.

\begin{center}
\fbox{\parbox{0.85\textwidth}{
Once we are told the total number of cars, each car can be treated like a trial:
\[
\text{northbound with probability } \frac{3}{5},
\qquad
\text{southbound with probability } \frac{2}{5}.
\]
}}
\end{center}

That is why the binomial distribution appears.

\section{Question: If Two Cars Pass in a Minute, What Is the Probability Both Are Northbound?}

We are given:

\[
\text{total number of cars in the minute}=2.
\]

Each of those 2 cars is northbound with probability
\[
\frac{3}{5}.
\]

So the number of northbound cars among the 2 total cars is

\[
X\sim \operatorname{Binomial}\left(2,\frac{3}{5}\right).
\]

We want both cars to be northbound:

\[
P(X=2).
\]

Using the binomial formula,

\[
P(X=k)=\binom{n}{k}p^k(1-p)^{n-k}.
\]

Here,

\[
n=2,\qquad k=2,\qquad p=\frac{3}{5}.
\]

Therefore,

\[
P(X=2)=\binom{2}{2}
\left(\frac{3}{5}\right)^2
\left(\frac{2}{5}\right)^0.
\]

Since

\[
\binom{2}{2}=1,
\qquad
\left(\frac{2}{5}\right)^0=1,
\]

we get

\[
P(X=2)=\left(\frac{3}{5}\right)^2.
\]

Thus,

\[
\boxed{P(\text{both northbound}\mid \text{two cars pass})=\frac{9}{25}.}
\]

\section{Meaning of the Binomial Value}

The binomial value is found from:

\[
\binom{n}{k}p^k(1-p)^{n-k}.
\]

\begin{itemize}
    \item \(n\) = total number of trials/events given.
    \item \(k\) = number of successes we want.
    \item \(p\) = probability of success for each event.
\end{itemize}

In this car example:

\begin{itemize}
    \item \(n=2\): two cars passed.
    \item \(k=2\): we want both to be northbound.
    \item \(p=3/5\): probability a given car is northbound.
\end{itemize}

\section{Question: If Four Cars Pass Between 11:55 and 12:05, What Is the Probability That Two Pass Before Noon and Two After Noon?}

The time interval from 11:55 to 12:05 has length 10 minutes.

It is split into two equal parts:

\[
11:55 \text{ to } 12:00 = 5 \text{ minutes},
\]

\[
12:00 \text{ to } 12:05 = 5 \text{ minutes}.
\]

We are given:

\[
\text{total number of cars in 10 minutes}=4.
\]

Since the two subintervals have equal length, each of the 4 cars is equally likely to fall before noon or after noon.

So

\[
P(\text{before noon})=\frac{5}{10}=\frac{1}{2},
\]

and

\[
P(\text{after noon})=\frac{5}{10}=\frac{1}{2}.
\]

Let

\[
Y=\text{number of cars before noon}.
\]

Then

\[
Y\sim \operatorname{Binomial}\left(4,\frac{1}{2}\right).
\]

We want exactly 2 before noon:

\[
P(Y=2).
\]

Using the binomial formula:

\[
P(Y=2)=\binom{4}{2}
\left(\frac{1}{2}\right)^2
\left(\frac{1}{2}\right)^2.
\]

Now,

\[
\binom{4}{2}=\frac{4!}{2!2!}=6.
\]

Therefore,

\[
P(Y=2)=6\left(\frac{1}{2}\right)^4
=\frac{6}{16}
=\frac{3}{8}.
\]

Thus,

\[
\boxed{P(\text{two before noon and two after noon}\mid \text{four cars pass})=\frac{3}{8}.}
\]

\section{Why This Is Binomial}

Once we are told that exactly 4 cars passed in the whole 10-minute interval, the actual Poisson count part is already fixed.

Now we only ask:

\begin{center}
How are these 4 cars distributed between the two parts of the interval?
\end{center}

Each car independently lands in the first 5-minute interval with probability \(1/2\), and in the second 5-minute interval with probability \(1/2\).

That is exactly a binomial setting.

\section{General Rule: Conditional on Total Count}

Suppose \(N(T)=n\) events occurred in a time interval of length \(T\).

If a subinterval has length \(a\), then each event has probability

\[
\frac{a}{T}
\]

of falling in that subinterval.

So the number of events in that subinterval, conditional on the total \(n\), follows

\[
\operatorname{Binomial}\left(n,\frac{a}{T}\right).
\]

\[
\boxed{
N(a)\mid N(T)=n
\sim
\operatorname{Binomial}\left(n,\frac{a}{T}\right).
}
\]

\section{Now About the 10th Event Question}

\textbf{Question.}
Compute the probability that the 10th event occurs 2 or more time units after the 9th event.

Let

\[
S_9=\text{time of the 9th event},
\]

\[
S_{10}=\text{time of the 10th event}.
\]

Then the time between the 9th and 10th event is

\[
T_{10}=S_{10}-S_9.
\]

The phrase:

\begin{center}
\textit{the 10th event occurs 2 or more time units after the 9th event}
\end{center}

means:

\[
S_{10}-S_9\geq 2.
\]

But

\[
T_{10}=S_{10}-S_9.
\]

Therefore, the event is

\[
T_{10}\geq 2.
\]

So

\[
\boxed{
P(\text{10th event occurs 2 or more time units after 9th})
=P(T_{10}\geq 2).
}
\]

\section{Why \(T_{10}\sim \operatorname{Exponential}(\lambda)\)}

In a Poisson process, the waiting times between consecutive events are independent exponential random variables.

That means:

\[
T_1=S_1-S_0,
\]

\[
T_2=S_2-S_1,
\]

\[
T_3=S_3-S_2,
\]

and so on all have the same distribution:

\[
T_i\sim \operatorname{Exponential}(\lambda).
\]

So the gap between the 9th and 10th event is also exponential:

\[
T_{10}=S_{10}-S_9\sim \operatorname{Exponential}(\lambda).
\]

\section{Calculating \(P(T_{10}\geq 2)\)}

For an exponential random variable,

\[
P(T\geq t)=e^{-\lambda t}.
\]

Here \(t=2\).

Therefore,

\[
P(T_{10}\geq 2)=e^{-\lambda(2)}.
\]

So

\[
\boxed{P(T_{10}\geq 2)=e^{-2\lambda}.}
\]

\section{Did We Really Use the Fact That It Is the 10th Event?}

Yes and no.

We used it to identify the gap:

\[
T_{10}=S_{10}-S_9.
\]

But after that, the calculation is the same for any consecutive pair of events.

For example:

\[
S_2-S_1,\quad S_5-S_4,\quad S_{10}-S_9
\]

all have the same exponential distribution.

This is because a Poisson process has independent and identically distributed interarrival times.

\section{Practice Questions With Answers}

\subsection*{Question 1}

Northbound cars pass at rate 4 per minute and southbound cars pass at rate 6 per minute. If 3 cars pass in a minute, what is the probability that exactly 2 are northbound?

\textbf{Answer.}

Total rate:

\[
4+6=10.
\]

Probability a given car is northbound:

\[
p=\frac{4}{10}=\frac{2}{5}.
\]

Given 3 total cars, let

\[
X=\text{number of northbound cars}.
\]

Then

\[
X\sim \operatorname{Binomial}\left(3,\frac{2}{5}\right).
\]

We want

\[
P(X=2).
\]

So

\[
P(X=2)=\binom{3}{2}
\left(\frac{2}{5}\right)^2
\left(\frac{3}{5}\right)^1.
\]

Thus,

\[
P(X=2)=3\cdot \frac{4}{25}\cdot \frac{3}{5}
=\frac{36}{125}.
\]

\[
\boxed{P(X=2)=\frac{36}{125}.}
\]

\subsection*{Question 2}

Six calls arrive between 1:00 and 1:30. What is the probability that exactly 2 arrived between 1:00 and 1:10?

\textbf{Answer.}

The full interval length is 30 minutes.

The subinterval length is 10 minutes.

So

\[
p=\frac{10}{30}=\frac{1}{3}.
\]

Given total calls:

\[
n=6.
\]

Let

\[
X=\text{number of calls in the first 10 minutes}.
\]

Then

\[
X\sim \operatorname{Binomial}\left(6,\frac{1}{3}\right).
\]

We want

\[
P(X=2).
\]

So

\[
P(X=2)=\binom{6}{2}
\left(\frac{1}{3}\right)^2
\left(\frac{2}{3}\right)^4.
\]

\[
\binom{6}{2}=15.
\]

Therefore,

\[
P(X=2)=15\cdot \frac{1}{9}\cdot \frac{16}{81}
=\frac{240}{729}
=\frac{80}{243}.
\]

\[
\boxed{P(X=2)=\frac{80}{243}.}
\]

\subsection*{Question 3}

Events occur according to a Poisson process with rate \(\lambda=3\). Find the probability that the 7th event occurs more than 4 time units after the 6th event.

\textbf{Answer.}

The gap between the 6th and 7th events is

\[
T_7=S_7-S_6.
\]

For a Poisson process,

\[
T_7\sim \operatorname{Exponential}(3).
\]

We want

\[
P(T_7>4).
\]

For exponential waiting times,

\[
P(T>t)=e^{-\lambda t}.
\]

So

\[
P(T_7>4)=e^{-3(4)}=e^{-12}.
\]

\[
\boxed{P(T_7>4)=e^{-12}.}
\]

\subsection*{Question 4}

Events occur according to a Poisson process with rate \(\lambda=2\). What is the probability that the time between the 20th and 21st events is at least 0.5 time units?

\textbf{Answer.}

The time between the 20th and 21st events is

\[
T_{21}=S_{21}-S_{20}.
\]

Since

\[
T_{21}\sim \operatorname{Exponential}(2),
\]

we want

\[
P(T_{21}\geq 0.5).
\]

Thus,

\[
P(T_{21}\geq 0.5)=e^{-2(0.5)}=e^{-1}.
\]

\[
\boxed{P(T_{21}\geq 0.5)=e^{-1}.}
\]

\section{Summary}

\begin{itemize}
    \item If the total number of events is given, splitting them into categories or subintervals often becomes binomial.
    \item For category splitting, use probability = category rate / total rate.
    \item For time splitting, use probability = subinterval length / total interval length.
    \item The time between consecutive Poisson events is exponential.
    \item ``10th event occurs 2 or more time units after the 9th event'' means \(S_{10}-S_9\geq 2\), or \(T_{10}\geq 2\).
\end{itemize}

\end{document}
